% Options for packages loaded elsewhere
\PassOptionsToPackage{unicode}{hyperref}
\PassOptionsToPackage{hyphens}{url}
\documentclass[
  a4paper,
]{ltjsarticle}
\usepackage{xcolor}
\usepackage[margin=25mm]{geometry}
\usepackage{amsmath,amssymb}
\setcounter{secnumdepth}{-\maxdimen} % remove section numbering
\usepackage{iftex}
\ifPDFTeX
  \usepackage[T1]{fontenc}
  \usepackage[utf8]{inputenc}
  \usepackage{textcomp} % provide euro and other symbols
\else % if luatex or xetex
  \usepackage{unicode-math} % this also loads fontspec
  \defaultfontfeatures{Scale=MatchLowercase}
  \defaultfontfeatures[\rmfamily]{Ligatures=TeX,Scale=1}
\fi
\usepackage{lmodern}
\ifPDFTeX\else
  % xetex/luatex font selection
\fi
% Use upquote if available, for straight quotes in verbatim environments
\IfFileExists{upquote.sty}{\usepackage{upquote}}{}
\IfFileExists{microtype.sty}{% use microtype if available
  \usepackage[]{microtype}
  \UseMicrotypeSet[protrusion]{basicmath} % disable protrusion for tt fonts
}{}
\makeatletter
\@ifundefined{KOMAClassName}{% if non-KOMA class
  \IfFileExists{parskip.sty}{%
    \usepackage{parskip}
  }{% else
    \setlength{\parindent}{0pt}
    \setlength{\parskip}{6pt plus 2pt minus 1pt}}
}{% if KOMA class
  \KOMAoptions{parskip=half}}
\makeatother
\usepackage{longtable,booktabs,array}
\newcounter{none} % for unnumbered tables
\usepackage{calc} % for calculating minipage widths
% Correct order of tables after \paragraph or \subparagraph
\usepackage{etoolbox}
\makeatletter
\patchcmd\longtable{\par}{\if@noskipsec\mbox{}\fi\par}{}{}
\makeatother
% Allow footnotes in longtable head/foot
\IfFileExists{footnotehyper.sty}{\usepackage{footnotehyper}}{\usepackage{footnote}}
\makesavenoteenv{longtable}
\setlength{\emergencystretch}{3em} % prevent overfull lines
\providecommand{\tightlist}{%
  \setlength{\itemsep}{0pt}\setlength{\parskip}{0pt}}
\usepackage{bookmark}
\IfFileExists{xurl.sty}{\usepackage{xurl}}{} % add URL line breaks if available
\urlstyle{same}
\hypersetup{
  hidelinks,
  pdfcreator={LaTeX via pandoc}}

\author{}
\date{}

\begin{document}

\section{無名二チャネル閉鎖波系におけるロック動力学}\label{ux7121ux540dux4e8cux30c1ux30e3ux30cdux30ebux9589ux9396ux6ce2ux7cfbux306bux304aux3051ux308bux30edux30c3ux30afux52d5ux529bux5b66}

\subsection{量子化機構（Arnold
tongue）の実在の実証と、位数選択機構の消去法------タング幅実測・固定点地図・観測選択の時計継承}\label{ux91cfux5b50ux5316ux6a5fux69cbarnold-tongueux306eux5b9fux5728ux306eux5b9fux8a3cux3068ux4f4dux6570ux9078ux629eux6a5fux69cbux306eux6d88ux53bbux6cd5ux30bfux30f3ux30b0ux5e45ux5b9fux6e2cux56faux5b9aux70b9ux5730ux56f3ux89b3ux6e2cux9078ux629eux306eux6642ux8a08ux7d99ux627f}

\textbf{著者:} 木原 範昭\\
\textbf{ORCID:} 0009-0004-6753-4020\\
\textbf{日付:} 2026年8月3日\\
\textbf{Version DOI:} \texttt{10.5281/zenodo.21764000}\\
\textbf{Concept DOI:} \texttt{10.5281/zenodo.21763999}\\
\textbf{位置づけ:} 「波の情報読出し」シリーズ・二チャネル交換散乱系
追加論文 v1（第9論文系列・\textbf{三部作の第三部}。第一部＝二文法分解
{[}A{]}、第二部＝数え上げ読出しの構造 {[}B{]}）

\begin{center}\rule{0.5\linewidth}{0.5pt}\end{center}

\subsection{要旨}\label{ux8981ux65e8}

\textbf{（零次 no-go）}
無名二チャネル閉鎖波系の標準読出し（位相盲目のビン別パワー比）は運動の厳密な定数であり（ドリフト
\(\le2.2\times10^{-16}\)、定理は {[}A{]}）、結合角 \(\theta\)
の動力学もロックも存在しない。

\textbf{（量子化機構の実在）} 位相感受読出し（対称チャネル \(X=A+B\)
のパワー比、保存的作業仮説として導入）の下では \(\theta\)
が発展し、\textbf{本物の Arnold tongue が実在する}：初期振幅の有限区間
\([2.518,\ \ge3.44]\)（幅
\(\ge0.92\)）にわたり、後期回転数の平均が厳密に \(\theta/\pi=1/3\)
にロックする（800衝突で機械精度、\(\theta\)
自身は振動する周期軌道＝有理平均回転）。根の規約で
\((m,n)=(1,6)\)、位数6。数え上げ構成による位数 3・4・6
のロック（{[}B{]}、結晶学的制限の完全メニュー）と合わせ、\textbf{量子化台地の機構は、結合を外部から与えない内生結合系に実在する}。

\textbf{（選択機構の消去法）}
しかしロックは\textbf{選択}の機構ではない。(i)
タング幅のモデル内実測：分母3の幅 \(\ge0.92\) に対し、探索帯で分母
\(\ge4\) のタングは幅 \(<0.002\)------\textbf{幅比
\(>461\)}。分母124の重みは無視できる。(ii)
タング下端直下は回転数が非単調に変動する超臨界領域であり、目標
\(\rho=39/124\) への最接近 \(3.3\times10^{-4}\) でプラトー非検出。(iii)
第二の動的候補・固定点型読出し（交差スペクトル）の吸引先は連続帯
\([0.016,0.450]\) に分布し\textbf{有理集中なし}（0/41）------消去。(iv)
観測選択（初期状態への回帰を確認する反復観測）は、厳密有限位数根で忠実度が厳密に1となり、\textbf{選択が観測時計の約数構造に継承される}ことを直接示す：時計
\(J=8\) の観測者には位数124の回帰が見えず（\(F_8=0.0026\)）、\(J=248\)
の観測者には見える（\(F_{248}=1.000000000000001\)）。ただし忠実度地形は根の近傍で滑らかであり、有限回の観測は回帰\textbf{盆地}までしか絞れない（\(n=20\)
で集中率 \(1.001\)、予言P2の反証として記録）。

\textbf{（帰結）}
量子化（なぜ飛び飛びか）はロックが説明し、選択（なぜこの値か）はロックでは説明できない。消去法の完了により、位数124の選択問題は\textbf{「観測時計はなぜ248と可換か」}という反証可能な単一の問いに変換された。この最終形は親論文（第9論文本体）が引き受ける。

\textbf{キーワード:} Arnold
tongue、モードロック、悪魔の階段、超臨界円写像、回帰忠実度、観測選択、量子再帰、再現可能計算

\begin{center}\rule{0.5\linewidth}{0.5pt}\end{center}

\subsection{0. 結論}\label{ux7d50ux8ad6}

\[
\boxed{
\begin{aligned}
&\text{位相盲目読出しの下で }\theta\text{ 動力学は存在しない（no-go、}2.2\times10^{-16}\text{）。}\\
&\text{位相感受読出しの下で Arnold tongue が実在：}\rho=1/3\text{ 厳密・幅}\ge0.92\text{（量子化の機構）。}\\
&\text{幅比}>461\text{・固定点消去・観測選択の時計継承により、ロックも固定点も位数124を選択しない。}\\
&\text{選択問題の最終形＝「観測時計はなぜ248と可換か」。}
\end{aligned}
}
\]

\subsection{本稿の位置づけ（三部構成）}\label{ux672cux7a3fux306eux4f4dux7f6eux3065ux3051ux4e09ux90e8ux69cbux6210}

第一部 {[}A{]}
は流れの二項恒等分解と二文法（加法/積・符号・中性化・階層）を、第二部
{[}B{]} は数え上げ読出しの構造（メタ勘定則・規約監査・Niven
交点・整数対読出し）を確立する。本稿（第三部）は動力学を扱う：量子化台地の機構の実在と、位数選択の候補機構の消去法である。三部は同一の再現パッケージを共有する。

\subsection{1. 研究課題}\label{ux7814ux7a76ux8ab2ux984c}

{[}B{]} により、数え上げ構成が到達できるロックは Niven の交点（位数 3,
4, 6）に限られる。一方、物理的な α に対応する根 \(R_{124,23}\)
は位数124であり
{[}2{]}、この選択は数え上げの外にある。第9論文系列の中心問題は「位数124は何が選ぶのか」であり、候補は
(a) 静的な状態構造（{[}B{]} で反証済）、(b) 動的ロック、(c)
固定点型アトラクタ、(d)
観測（回帰確認）による選択、の四つに整理される。本稿は (b)(c)(d)
を実験で判定する。

\subsection{2. 系と設計境界}\label{ux7cfbux3068ux8a2dux8a08ux5883ux754c}

系・状態構成・標準読出し・回転は {[}A{]}
第2節と同一である（散乱本体・読出し無変更、目標値は状態構成のみに使用）。位相感受読出しの導入は保存的作業仮説として明示し（第4節）、観測選択の「観測」は初期状態対への回帰忠実度

\[
F_J=\frac{|\langle a_0|a_J\rangle+\langle b_0|b_J\rangle|^2}{(\langle a_0|a_0\rangle+\langle b_0|b_0\rangle)^2}
\]

の評価として定義する（力学への逆作用なし・純粋な読出し）。

\subsection{3. 零次 no-go}\label{ux96f6ux6b21-no-go}

標準読出しはビン別回転不変量 \(|A_k|^2+|B_k|^2\)
のみの関数であり、実直交回転の下で厳密に保存される（定理は {[}A{]}
第4節）。数値確認：400衝突×4振幅で \(\theta\) の最大ドリフト
\(2.2\times10^{-16}\)。ゆえに \(\theta\)
の動力学・ロック・選択は零次に存在しない。以降の全動力学は、この no-go
を破る読出しの明示的選択の上に立つ。

\subsection{4.
読出し分類（動力学面）}\label{ux8aadux51faux3057ux5206ux985eux52d5ux529bux5b66ux9762}

\(\theta\)
読出し7種を「自身が駆動する動力学の下での挙動」で分類した（{[}A{]}
第4節の表の動力学的側面）。保存は対角和のみ。動力学級のうち、交差スペクトル読出し（R4）は\textbf{固定点収束型}（初期から大きく動いた後に凍結。当初これを保存と誤認し、分類実験で訂正------付記）、\(X\)
パワー読出し（V1）は\textbf{振動・ロック型}である。以降、V1
をロック探査に、R4 を固定点探査に用いる。

\subsection{5. Arnold tongue
の実在（量子化機構）}\label{arnold-tongue-ux306eux5b9fux5728ux91cfux5b50ux5316ux6a5fux69cb}

\subsubsection{5.1
発見と精密化}\label{ux767aux898bux3068ux7cbeux5bc6ux5316}

V1
の下で初期B振幅を掃引すると、後期回転数の末尾平均が厳密に有理数へ張り付く有限区間が存在する：

\begin{itemize}
\tightlist
\item
  粗スキャン（61点×400衝突）でプラトー検出、精密スキャン（13点×800衝突・tail
  300）で確定：\textbf{振幅 2.6--3.3 で \(\overline{\theta/\pi}=1/3\)
  が機械精度（距離 0.0、最大 \(5.6\times10^{-17}\)）}
\item
  下端は 39/124 マイクロズーム（27点×1500衝突）で 2.518
  と確定、上端は追加スキャンで \(\ge3.44\)------\textbf{幅 \(\ge0.92\)}
\item
  \(\theta\) 自身は振動（末尾標準偏差
  \(\sim7\times10^{-2}\)）し、平均回転数のみが有理ロックする周期軌道------円写像のモードロックの標準的挙動
\end{itemize}

根の規約 \(\theta=\pi/2-\pi m/n\) で \((m,n)=(1,6)\)。{[}B{]}
の数え上げロック（位数 3, 4,
6）と合わせ、内生結合系における量子化台地の機構が確立した。

\subsubsection{5.2
タング幅のモデル内実測（選択への第一の否定）}\label{ux30bfux30f3ux30b0ux5e45ux306eux30e2ux30c7ux30ebux5185ux5b9fux6e2cux9078ux629eux3078ux306eux7b2cux4e00ux306eux5426ux5b9a}

分母 \(\ge4\) のタングは、探索帯 \([2.470, 2.522]\)（分解能
0.002・1500衝突）で一つも検出されなかった。したがってモデル内実測として

\[
\frac{\mathrm{width}(q=3)}{\mathrm{width}(q\ge4)}>\frac{0.92}{0.002}=461
\]

であり、幅は分母とともに急崩壊する。\textbf{分母124のタング重みは、外部理論（円写像の一般論）を引くまでもなく、自前のデータで無視できる}。

\subsubsection{5.3 超臨界診断と 39/124
上界（第二の否定）}\label{ux8d85ux81e8ux754cux8a3aux65adux3068-39124-ux4e0aux754cux7b2cux4e8cux306eux5426ux5b9a}

1/3 タング下端の直下では、回転数が隣接振幅点間 \(\sim10^{-2}\)
で非単調に変動する------タングが重なる超臨界領域の症状である。目標
\(\rho=39/124=0.3145161\ldots\) への最接近は \(3.3\times10^{-4}\)（振幅
2.504、1500衝突）で、プラトーは検出されない。素朴掃引による高位数タング探索はこの読出しでは原理的に不適であり、当初設計した結合補間キャンペーンは、5.2
の実測比により\textbf{結論に不要}となった（経緯を付記に記録）。

\subsection{6.
固定点地図（第三の候補の消去）}\label{ux56faux5b9aux70b9ux5730ux56f3ux7b2cux4e09ux306eux5019ux88dcux306eux6d88ux53bb}

R4 読出しの自己駆動下の吸引先 \(\theta^*\)
を初期振幅41点で測った（400衝突）。吸引先は連続帯
\(\theta^*/\pi\in[0.016, 0.450]\)
に分布し、\textbf{有理角への厳密集中は皆無}（0/41）。固定点型アトラクタは特別な値を選ばない------第三の候補は消去された。

\subsection{7.
観測選択（第四の候補：機構は実在、選択は時計に継承）}\label{ux89b3ux6e2cux9078ux629eux7b2cux56dbux306eux5019ux88dcux6a5fux69cbux306fux5b9fux5728ux9078ux629eux306fux6642ux8a08ux306bux7d99ux627f}

\subsubsection{7.1 設計と予言}\label{ux8a2dux8a08ux3068ux4e88ux8a00}

状態族（振幅族61点＋厳密根状態2点：\(R=1/2\) と
\(R_{124,23}\)、逆算探索で構成）に対し、時計 \(J\in\{8,248\}\)
での回帰忠実度 \(F_J\) と、\(n=20\) 回の反復観測後の生存重み \(F_J^n\)
を測る。予言（測定前固定）：P1＝厳密根でのみ \(F=1\)（\(R=1/2\) は
\(8\mid248\) により両時計で、124根は \(J=248\) のみで）、P2＝\(n=20\)
で回帰状態への集中率 \(>10^3\)。

\subsubsection{7.2 結果}\label{ux7d50ux679c}

\textbf{P1
は全成立}：\(F_8(R{=}1/2)=0.999999999999999\)、\(F_{248}(\text{124根})=1.000000000000001\)、そして\textbf{時計選択性}------\(F_8(\text{124根})=0.0026\)。すなわち\textbf{どの回帰が「見える」かは観測時計の約数構造が決める}：\(J=8\)
の観測者は位数124の回帰に原理的に盲目であり、\(J=248\)
の観測者は位数8と124の両方を選ぶ。選択は状態でなく観測時計に継承される。

\textbf{P2 は反証}（設計どおり記録）：\(n=20\) での集中率は
\(1.001\)。忠実度地形は根の近傍で滑らかであり、近傍状態も
\(F\approx1-\varepsilon\)
で生き残るため、\textbf{有限回の観測は回帰盆地を選べても厳密点は漸近的にしか選べない}。これは公開論文
{[}2{]}
の「深さは理想演算においてのみ無限大」という記述と整合する、観測選択の定量的限界である。

\subsubsection{7.3
帰結：選択問題の座標変換}\label{ux5e30ux7d50ux9078ux629eux554fux984cux306eux5ea7ux6a19ux5909ux63db}

四候補の判定が完了した：静的構造（×・{[}B{]}）、動的ロック（×・幅比461）、固定点（×・連続帯）、観測選択（機構は実在するが、選択は観測時計の可換性
\(P\mid J\) に継承される）。ゆえに「なぜ位数124か」は

\[
\boxed{\text{「観測時計はなぜ 248 と可換か」}}
\]

という単一の問いに変換された。これは後退ではなく前進である：問いが状態空間の探索から、観測者（レジスタ・分解能）の構造の問題へ移った。この最終形は親論文（第9論文本体：フェルミオンの生成構造・波束収縮）が引き受ける。

\subsection{8. 主張}\label{ux4e3bux5f35}

\textbf{主張1.}
位相盲目読出しの下で結合角の動力学は存在しない（no-go・定理＋数値
\(2.2\times10^{-16}\)）。

\textbf{主張2.} 位相感受読出しの下で Arnold tongue
が実在する：\(\rho=1/3\) 厳密・幅
\(\ge0.92\)・\((m,n)=(1,6)\)。数え上げロック（位数3,4,6
{[}B{]}）と合わせ、\textbf{量子化（なぜ飛び飛びか）の機構は内生結合系に実在する}。

\textbf{主張3.}
ロックは選択（なぜこの値か）の機構ではない：タング幅はモデル内実測で分母3→分母\$\ge\$4
の間に \(>461\) 倍崩壊し、位数124の重みは無視できる。

\textbf{主張4.} 固定点型チャネルは選択しない（吸引先連続帯・有理集中
0/41）。

\textbf{主張5.}
観測選択は機構として実在し（厳密根で忠実度1）、選択を\textbf{観測時計の約数構造に継承する}（時計選択性の直接実証）。ただし有限回観測は回帰盆地までしか絞れない（P2反証）。ゆえに選択問題の最終形は「観測時計はなぜ248と可換か」である。

\subsection{9.
既知系との関係}\label{ux65e2ux77e5ux7cfbux3068ux306eux95a2ux4fc2}

モードロック・悪魔の階段・タング幅の分母依存は円写像の確立された数学であり
{[}13{]}、位相ロックによる量子化台地は超伝導 Shapiro 段差 {[}12{]}
として物理に実在する。本稿の寄与は、(i)
結合が装置定数でなく状態から読まれる内生結合系でロックを実証したこと、(ii)
タング幅・固定点・観測選択を同一系で測り選択候補の消去法を完了したこと、(iii)
観測時計の約数構造への選択継承を直接示したこと、である。

\subsection{10. 未解決問題}\label{ux672aux89e3ux6c7aux554fux984c}

\begin{enumerate}
\def\labelenumi{\arabic{enumi}.}
\tightlist
\item
  \textbf{観測時計の可換性}：「なぜ248か」の最終形。観測者自身を系内の閉包（レジスタ）としてモデル化し、その固有周期が支配的な回帰周期と可換になる機構の探索------親論文へ
\item
  \textbf{V1
  読出しの導出}：位相感受読出しは作業仮説である。保存読出し分類（{[}A{]}
  第4節）は候補空間を狭めたが、力学が選ぶ読出しの第一原理はない
\item
  \textbf{亜臨界領域の階段}：結合補間 \(P_\lambda\)
  による完全階段の測定は、結論には不要となったが（5.2）、タング幅スケーリングの定量理論としては残る
\end{enumerate}

以上はいずれもモデル内部の未解決問題である。観測物理への接続として残る課題の全体像は、これらとは別系列として次節に一括整理する。

\subsection{11. 三部作を越える物理接続課題------三方向から 3+1
次元時空へ、二文法から重力・クーロン力へ}\label{ux4e09ux90e8ux4f5cux3092ux8d8aux3048ux308bux7269ux7406ux63a5ux7d9aux8ab2ux984cux4e09ux65b9ux5411ux304bux3089-31-ux6b21ux5143ux6642ux7a7aux3078ux4e8cux6587ux6cd5ux304bux3089ux91cdux529bux30afux30fcux30edux30f3ux529bux3078}

本稿は三部作の最終稿であるから、モデル内部で確立された範囲と、観測物理への接続として残る範囲を、最後に一枚の課題表として明示する。

まず論理階層を区別する。前節までに残された\textbf{内部選択問題}（「観測時計はなぜ248と可換か」）は、モデル内部で特定の有限位数根が選択される理由を問う。本節の\textbf{物理接続問題}は、そのように選択された内部構造が、観測される
3+1
次元時空・重力・クーロン力としてどのように実現されるかを問う。両者は競合しない別系列の課題である。

導出済みのものを再発見する必要はない：三方向の自発的発生（本系列前段・親論文が報告する準安定化過程）、調和閉鎖による逆二乗則
{[}3{]}、重力型・電荷型二文法の一意な分化（{[}A{]}
6.5節、反転割当ての排除）、数え上げ・有限位数・ロック・観測時計の構造（{[}B{]}・本稿）。残るのは、これらの内部構造を観測物理へ写す\textbf{実現写像}の構成であり、具体的には以下の課題である。

\textbf{（1）三方向の局所空間化.} 準安定時に立ち上がる三方向の生成作用を
\(T_1,T_2,T_3\)
とし、それらが三つの内部モードにとどまらず局所三次元空間の基底を成すことを示す。零次候補は
Gram 計量 \(g_{ij}(\Psi)=\mathrm{Re}\langle T_i\Psi|T_j\Psi\rangle\)
であり、準安定状態族に対して
\(\mathrm{rank}\,g=3\)・正定値性が成立するかを検証する。三方向の番号付けに物理的意味を与えてはならず（無名性）、基底回転
\(SO(3)\)
に対して同値な空間として定義できることが要件である。\textbf{三方向の存在と三次元空間の成立は同じことではない}------これが本節全体の要検査の核である。

\textbf{（2）局所三方向の大域化と可積分性.}
各閉包・各局所状態で得られる三方向が共通の空間を形成するか。\([T_i,T_j]\)
の構造、局所基底の平行移動と経路依存性（ホロノミー）、Frobenius
型の可積分条件を調べ、局所三方向が一つの大域的空間へ接続可能かを判定する。

\textbf{（3）時間方向と 3+1 次元時空の構成.}
本系には衝突回数・位相進行・有限位数回帰・観測時計 \(J\)
という複数の時間候補がある。何が（またはどの組合せが）物理的時間を担うかを確定し、三つの正定値方向と合わせて
Lorentz 型計量
\(ds^2=-c_{\mathrm{eff}}^2 d\tau^2+g_{ij}dx^i dx^j\)、因果構造、有限伝播速度が内生的に生じるかを検証する。

\textbf{（4）距離の同一性.}
調和閉鎖の距離・三方向から構成される幾何距離・観測者が読む距離の同一性

\[
r_{\mathrm{harmonic}}=r_{\mathrm{geometric}}=r_{\mathrm{observed}}
\]

を導出する（{[}A{]} 未解決2の \(\Delta\theta\leftrightarrow r\)
辞書の完成形）。スカラー \(r\) だけでは足りない：二閉包を結ぶ方向
\(\hat{\mathbf r}\)
と勾配演算を構成し、\(\nabla(1/r)=-\hat{\mathbf r}/r^2\)
に相当する\textbf{方向付き}逆二乗場が創発三方向上で定義できることを示す必要がある。

\textbf{（5）重力型文法から物理的重力への写像.}
大きさの項の位相盲目・加法的源読出しを物理的な質量・エネルギー源へ写し、源が創発計量・接続をどう変形するか（大きさ型源
→ \(\delta g_{\mu\nu}\) または \(\delta\Gamma^\mu{}_{\nu\rho}\) →
閉包の運動）を定式化する。最初の検証目標は、弱場・低速極限の Newton
型加速度、試験閉包の内部構造に依存しない普遍的落下、慣性側と重力側の読出しの一致、および結合とスケール比
\(R/R_0\) を結ぶ階層指数の内生的導出（{[}A{]} のスケール辞書）である。

\textbf{（6）電荷型文法から物理的クーロン力への写像.} 重なりの項
\(\mathrm{Re}\langle a|b\rangle\sim|q_A||q_B|\cos(\phi_B-\phi_A)\)
を電荷型源頂点として読み、調和閉鎖の媒介 {[}3{]}
と受け手側の応答に接続して、創発空間上で
\(\mathbf F_{AB}=K\,q_Aq_B\,\hat{\mathbf r}_{AB}/r_{AB}^2\)
を導出する（源頂点 × 媒介 ×
応答頂点の構造）。個別に確定すべき項目：相対位相の符号が接近・離反のどちらへ写るか（Bjerknes
型か電磁気型か------{[}A{]}
の自動判定予言）、チャネル非共有の中性化が物理的電気中性へ写るか、電荷型源の保存則と
Gauss 型流束則、\(1-R_\ast=\sin^2(23\pi/124)\)
が単位規約込みで観測素電荷へ写る経路、時間依存位相の導入で静的クーロン場を越えて磁気的・放射的成分が生じるか。

\textbf{（7）二相互作用に共通する時空辞書.}
重力型と電荷型に別々の距離・別々の空間を割り当ててはならない。同一の創発計量・同一の
\(r\)・同一の \(\tau\)
の上で、大きさ型文法→重力応答と重なり型文法→電荷応答が\textbf{同時に}成立し、かつこの対応が三方向の基底回転・A/B
チャネル交換・倍音の名称変更・読出し規約変更・観測者の座標選択に対して不変または共変であることを示す。

\textbf{（8）物理接続写像の反証条件.}
接続は類似性では認定しない。以下のいずれかが成立すれば、物理的重力・クーロン力との同定を棄却する：三方向から
rank 3
の安定計量が構成できない／局所三方向が大域的に可積分でない／調和閉鎖の距離と創発幾何距離が一致しない／大きさ型源が普遍的な幾何応答を生まない／重なり型の位相符号が一貫した引力・斥力へ写らない／同じ電荷量が源側と応答側の双方に現れない／二文法が同一の
3+1 次元時空上で同時に作用しない。

以上は本三部作の内部成果を未確定へ戻す課題ではなく、導出済みの前幾何学的構造を観測量へ写す次段の接続課題である。本系列の現在地は次のとおり整理される：

\[
\boxed{
\begin{aligned}
&\text{三方向の自発発生} && \text{導出済み（本系列前段）}\\
&\text{逆二乗媒介} && \text{導出済み [3]}\\
&\text{重力型・電荷型の二文法（割当て一意）} && \text{導出・検証済み [A]}\\
&\text{数え上げ・有理量子化・観測時計の構造} && \text{導出・限界検証済み（[B]・本稿）}\\
&\text{三方向}\to 3+1\text{ 次元時空} && \text{未接続（課題 1–4）}\\
&\text{二文法}\to\text{物理的重力・クーロン力} && \text{未接続（課題 4–7）}
\end{aligned}
}
\]

三部作の結論は三段で閉じる：\textbf{二文法は自発的に分化し、割当ては一意である（{[}A{]}）。その読出しは数え上げ・有限位数・ロック・観測時計の構造を持つ（{[}B{]}・本稿）。物理時空と物理的力への実現写像は、上記
8 課題として次段------親論文とその後続------が引き受ける。}

\subsection{12. 再現性}\label{ux518dux73feux6027}

全実験のランナー・全測定CSV・図はコミット済みであり、{[}A{]}
と同一の規律（SHA検証付き中核参照・予言の測定前固定・アンカー再現・反証と訂正の保存）に従う。

{\def\LTcaptype{none} % do not increment counter
\begin{longtable}[]{@{}
  >{\raggedright\arraybackslash}p{(\linewidth - 4\tabcolsep) * \real{0.3333}}
  >{\raggedright\arraybackslash}p{(\linewidth - 4\tabcolsep) * \real{0.3333}}
  >{\raggedright\arraybackslash}p{(\linewidth - 4\tabcolsep) * \real{0.3333}}@{}}
\toprule\noalign{}
\begin{minipage}[b]{\linewidth}\raggedright
実験
\end{minipage} & \begin{minipage}[b]{\linewidth}\raggedright
フォルダ
\end{minipage} & \begin{minipage}[b]{\linewidth}\raggedright
コミット
\end{minipage} \\
\midrule\noalign{}
\endhead
\bottomrule\noalign{}
\endlastfoot
モードロック探査（Z₆発見・no-go） &
\texttt{mode\_locking\_probe\_pre\_v1} & 56453653 \\
1/3タング精密スキャン & 同上
\texttt{run\_one\_third\_tongue\_fine\_scan\_v1} & 15ee20cb \\
読出し分類＋階段＋39/124ズーム & \texttt{tongue\_spectroscopy\_pre\_v1}
& 913756fe, 15ee20cb \\
タング幅実測（上端・幅比461） & \texttt{tongue\_width\_measurement\_v1}
& 47e2da02 \\
固定点地図 & \texttt{partial\_sharing\_fixed\_point\_v1}（E7部） &
47e2da02 \\
観測選択（時計選択性） & \texttt{observation\_selection\_v1} &
47e2da02 \\
\end{longtable}
}

\begin{center}\rule{0.5\linewidth}{0.5pt}\end{center}

\section{参考文献}\label{ux53c2ux8003ux6587ux732e}

\subsection{自己引用}\label{ux81eaux5df1ux5f15ux7528}

\begin{enumerate}
\def\labelenumi{\arabic{enumi}.}
\tightlist
\item
  木原範昭「無名等振幅複合波モデル基本公理系 v6」Version DOI:
  \texttt{10.5281/zenodo.21465984}, Concept DOI:
  \texttt{10.5281/zenodo.21315735}, 2026.
\item
  木原範昭「反復交換散乱における有限位数共鳴の発見」Version DOI:
  \texttt{10.5281/zenodo.21421367}, Concept DOI:
  \texttt{10.5281/zenodo.21421366}, 2026.
\item
  木原範昭「AB二体閉鎖位相系における未来位相位置加速度写像と調和閉鎖による逆二乗則
  v4」Version DOI: \texttt{10.5281/zenodo.21468270}, Concept DOI:
  \texttt{10.5281/zenodo.21441081}, 2026. （零次
  no-go→保存的作業仮説→独立検証の三段手順の先例）
\item
  木原範昭「無名二チャネル閉鎖波系における相互作用の二文法分解」（三部作第一部）Version
  DOI: \texttt{10.5281/zenodo.21763996}, Concept DOI:
  \texttt{10.5281/zenodo.21763995}, 2026.
\item
  木原範昭「無名二チャネル閉鎖波系における数え上げ読出しの構造」（三部作第二部）Version
  DOI: \texttt{10.5281/zenodo.21763998}, Concept DOI:
  \texttt{10.5281/zenodo.21763997}, 2026.
\item
  木原範昭「交換干渉散乱行列\ldots 局在性交換予備実験総括 v1」Version
  DOI: \texttt{10.5281/zenodo.21333768}, 2026. （エンジンのコード出自）
\item
  木原範昭「交換散乱係数R集中と微細構造定数対応候補の数値実験
  v1」Version DOI: \texttt{10.5281/zenodo.21396761}, 2026. （System A
  コード出自）
\end{enumerate}

\subsection{外部参考文献}\label{ux5916ux90e8ux53c2ux8003ux6587ux732e}

\begin{enumerate}
\def\labelenumi{\arabic{enumi}.}
\setcounter{enumi}{11}
\tightlist
\item
  S. Shapiro, \emph{Phys. Rev.~Lett.} \textbf{11}, 80 (1963). DOI:
  \texttt{10.1103/PhysRevLett.11.80}.
\item
  M. H. Jensen, P. Bak, and T. Bohr, \emph{Phys. Rev.~Lett.}
  \textbf{50}, 1637 (1983). DOI: \texttt{10.1103/PhysRevLett.50.1637}.
\end{enumerate}

\begin{center}\rule{0.5\linewidth}{0.5pt}\end{center}

\textbf{付記（訂正と経緯の記録）} (i) 交差スペクトル読出し R4
を「第二の保存読出し」と誤認→分類実験で固定点収束型と判明し訂正。(ii)
観測選択の集中予言 P2（\(n=20\) で
\(>10^3\)）は反証された------滑らかな忠実度地形による有限回観測の原理的限界として本文に記録。(iii)
当初設計した結合補間（\(P_\lambda\)）による亜臨界階段キャンペーンは、タング幅のモデル内実測（幅比
\(>461\)）が同じ結論をより直接に与えたため、結論に不要となった。判断の経緯ごと記録する。

\end{document}
